\lecture{36}{Градиентный трекинг}

\sourcevideo
    {https://www.youtube.com/playlist?list=PLOzRYVm0a65fhKDS8cYIC7YCuVGJ4FOJx}
    {Week 10: Lecture 36}

Распределённый градиентный спуск сочетает смешивание локальных
переменных с шагом по локальному градиенту, но при постоянном шаге в
общем случае сохраняет смещение. Метод градиентного трекинга вводит
дополнительную переменную, которая отслеживает средний градиент сети и
позволяет получить точную сходимость при постоянном шаге.

\section{Постановка и смещение DGD}

Пусть \(N\) агентов решают задачу
\[
    \min_{x\in\mathbb R^d} F(x),
    \qquad
    F(x)=\sum_{i=1}^{N}f_i(x),
\]
где агент \(i\) знает только функцию
\(f_i\colon\mathbb R^d\to\mathbb R\) и обменивается сообщениями с
соседями. Каждый агент хранит локальную оценку
\(x_i^k\in\mathbb R^d\), и требуется обеспечить
\[
    x_i^k\longrightarrow x^\star,
    \qquad
    x^\star\in\operatorname*{arg\,min}_{x}F(x).
\]

Распределённый градиентный спуск имеет вид
\[
    x_i^{k+1}
    =
    \sum_{j=1}^{N}w_{ij}x_j^k
    -\eta\nabla f_i(x_i^k).
\]
Первое слагаемое уменьшает рассогласование, второе --- локальную
функцию. При постоянном шаге эти действия конфликтуют: в глобальном
оптимуме обычно
\[
    \nabla f_i(x^\star)\ne0,
\]
хотя
\[
    \sum_{i=1}^{N}\nabla f_i(x^\star)=0.
\]
Поэтому локальные градиенты продолжают раздвигать агентов даже около
решения, и DGD может сходиться лишь к его окрестности.

\section{Алгоритм градиентного трекинга}

Каждый агент дополнительно хранит \(y_i^k\in\mathbb R^d\), которое
должно приближать средний текущий градиент сети. Пусть
\(W=(w_{ij})\in\mathbb R^{N\times N}\) --- примитивная дважды
стохастическая матрица смешивания:
\[
    W\mathbf1=\mathbf1,
    \qquad
    \mathbf1^{\mathsf T}W=\mathbf1^{\mathsf T}.
\]
При постоянном шаге \(\eta>0\) метод задаётся рекурсиями
\[
    x_i^{k+1}
    =
    \sum_{j=1}^{N}w_{ij}x_j^k
    -\eta y_i^k,
\]
\[
\begin{aligned}
    y_i^{k+1}
    ={}&
    \sum_{j=1}^{N}w_{ij}y_j^k
    \\
    &+
    \nabla f_i(x_i^{k+1})
    -
    \nabla f_i(x_i^k),
\end{aligned}
\]
с инициализацией
\[
    y_i^0=\nabla f_i(x_i^0).
\]

Обновление \(x_i^k\) использует не локальный градиент, а его
отслеживаемую сетевую оценку. В обновлении \(y_i^k\) смешивание
сочетается с приращением локального градиента. Без разности
\[
    \nabla f_i(x_i^{k+1})-\nabla f_i(x_i^k)
\]
обычный консенсус усреднял бы только начальные градиенты и не учитывал
бы движение точек \(x_i^k\).

На каждой итерации агент передаёт соседям два вектора: \(x_i^k\) и
\(y_i^k\). Наличие двух связанных рекурсий не делает метод методом
второго порядка: в нём нет ни гессианов, ни обратных гессианов, а
переменная \(y_i^k\) не является классическим импульсом.

\section{Инвариант и средняя динамика}

Главное тождество метода имеет вид
\[
    \sum_{i=1}^{N}y_i^k
    =
    \sum_{i=1}^{N}\nabla f_i(x_i^k)
\]
для всех \(k\). При \(k=0\) оно следует из инициализации. Если
равенство выполнено на итерации \(k\), то
\[
\begin{aligned}
    \sum_{i=1}^{N}y_i^{k+1}
    ={}&
    \sum_{i=1}^{N}\sum_{j=1}^{N}w_{ij}y_j^k
    \\
    &+
    \sum_{i=1}^{N}
    \bigl(
        \nabla f_i(x_i^{k+1})
        -
        \nabla f_i(x_i^k)
    \bigr).
\end{aligned}
\]
Столбцовая стохастичность даёт
\[
    \sum_{i=1}^{N}\sum_{j=1}^{N}w_{ij}y_j^k
    =
    \sum_{j=1}^{N}y_j^k,
\]
после чего индукционное предположение приводит к тому же тождеству на
итерации \(k+1\).

Определим средние
\[
    \overline x^k
    =
    \frac1N\sum_{i=1}^{N}x_i^k,
    \qquad
    \overline y^k
    =
    \frac1N\sum_{i=1}^{N}y_i^k.
\]
Тогда
\[
    \overline y^k
    =
    \frac1N
    \sum_{i=1}^{N}\nabla f_i(x_i^k)
\]
и, после усреднения первого обновления,
\[
    \overline x^{k+1}
    =
    \overline x^k-\eta\overline y^k
    =
    \overline x^k
    -
    \frac{\eta}{N}
    \sum_{i=1}^{N}\nabla f_i(x_i^k).
\]
Если локальные точки близки к среднему, то
\[
    \overline y^k
    \approx
    \frac1N\nabla F(\overline x^k),
\]
поэтому средняя переменная движется приблизительно как
централизованный градиентный спуск для функции \(F/N\).

\section{Матричная форма и рассогласование}

Объединим локальные переменные в столбцы
\[
    \mathbf x^k
    =
    \begin{pmatrix}
        x_1^k\\
        \vdots\\
        x_N^k
    \end{pmatrix},
    \qquad
    \mathbf y^k
    =
    \begin{pmatrix}
        y_1^k\\
        \vdots\\
        y_N^k
    \end{pmatrix},
\]
и положим
\[
    \mathbf g(\mathbf x)
    =
    \begin{pmatrix}
        \nabla f_1(x_1)\\
        \vdots\\
        \nabla f_N(x_N)
    \end{pmatrix}.
\]
Тогда
\[
    \mathbf x^{k+1}
    =
    (W\otimes I_d)\mathbf x^k
    -\eta\mathbf y^k,
\]
\[
\begin{aligned}
    \mathbf y^{k+1}
    ={}&
    (W\otimes I_d)\mathbf y^k
    \\
    &+
    \mathbf g(\mathbf x^{k+1})
    -
    \mathbf g(\mathbf x^k).
\end{aligned}
\]

Пусть
\[
    J=\frac1N\mathbf1\mathbf1^{\mathsf T}.
\]
Ошибки консенсуса равны
\[
    \mathbf e_x^k
    =
    \mathbf x^k-
    \mathbf1\otimes\overline x^k,
    \qquad
    \mathbf e_y^k
    =
    \mathbf y^k-
    \mathbf1\otimes\overline y^k.
\]
Матрица \(W-J\) действует только на неконсенсусные компоненты. Для
симметричной примитивной дважды стохастической матрицы
\[
    \rho=\lVert W-J\rVert_2<1,
\]
поэтому смешивание сжимает обе ошибки. В динамике
\(\mathbf e_y^k\) возникает возмущение, связанное с изменением
градиентов, но при их липшицевости оно контролируется изменением
\(\mathbf x^k\).

\section{Предельная точка и сходимость}

Предположим, что последовательности сходятся и локальные переменные
достигают консенсуса:
\[
    x_i^k\to x^\infty,
    \qquad
    y_i^k\to y_i^\infty.
\]
В пределе первое обновление даёт
\[
    \mathbf x^\infty
    =
    (W\otimes I_d)\mathbf x^\infty
    -\eta\mathbf y^\infty.
\]
На консенсусном подпространстве
\[
    (W\otimes I_d)\mathbf x^\infty=\mathbf x^\infty,
\]
поэтому
\[
    \mathbf y^\infty=0.
\]
Инвариант суммы отсюда влечёт
\[
    0
    =
    \sum_{i=1}^{N}\nabla f_i(x^\infty)
    =
    \nabla F(x^\infty).
\]
Для выпуклой функции \(F\) это условие означает глобальную
оптимальность. Именно здесь исчезает смещение DGD: в глобальном
оптимуме отслеживаемый средний градиент равен нулю, хотя отдельные
локальные градиенты могут оставаться ненулевыми.

\begin{theorem}[Линейная сходимость]
Пусть граф связен и неориентирован, \(W\) примитивна и дважды
стохастична, локальные градиенты липшицевы, глобальная функция \(F\)
сильно выпукла, а постоянный шаг \(\eta\) достаточно мал. Тогда
существуют \(C>0\) и \(q\in(0,1)\), для которых
\[
    \lVert x_i^k-x^\star\rVert
    \le
    Cq^k
\]
для всех агентов \(i\).
\end{theorem}

Допустимый шаг зависит от гладкости локальных функций и качества
смешивания, характеризуемого величиной \(\lVert W-J\rVert_2\).
Слишком большой шаг может вызвать колебания или расходимость даже при
сильной выпуклости.

\section{Пример, реализация и стоимость}

Для функций
\[
    f_i(x)=\frac12(x-i)^2,
    \qquad
    i=1,\ldots,5,
\]
имеем \(\nabla f_i(x)=x-i\) и
\[
    x^\star
    =
    \frac{1+2+3+4+5}{5}
    =3.
\]
При произвольных \(x_i^0\) и инициализации
\[
    y_i^0=x_i^0-i
\]
рекурсии принимают вид
\[
    x_i^{k+1}
    =
    \sum_jw_{ij}x_j^k-
    \eta y_i^k,
\]
\[
    y_i^{k+1}
    =
    \sum_jw_{ij}y_j^k
    +x_i^{k+1}-x_i^k.
\]
При подходящем шаге все локальные точки сходятся к \(3\).

Для неориентированного графа можно использовать веса Метрополиса:
\[
    w_{ij}
    =
    \frac1{1+\max\{d_i,d_j\}}
\]
для соседних \(i\ne j\) и
\[
    w_{ii}
    =
    1-
    \sum_{j\in\mathcal N_i}w_{ij}.
\]
Полученная матрица симметрична и дважды стохастична. В NumPy
смешивание записывается как \texttt{W @ x}; выражение
\texttt{W * x} выполняет покомпонентное умножение и не реализует
матричный шаг.

По каждому ребру на итерации передаются два вектора размерности \(d\),
поэтому суммарный объём обмена имеет порядок
\[
    O(|\mathcal E|d),
\]
но с примерно вдвое большей постоянной, чем у DGD. Локально агент
вычисляет один новый градиент и два смешивания; предыдущий градиент
можно хранить.

Базовая модель предполагает синхронные итерации, фиксированный граф и
точную передачу сообщений. Для ориентированных графов требуются
push-sum или другие корректирующие механизмы. При стохастических
градиентах переменная трекинга отслеживает шумную величину, поэтому в
анализе появляется дополнительная дисперсия.
